\section{polylog Indep}
%%% generated by the blueprint pipeline from TCSlib.BooleanAnalysis.polylogIndep
\subsection*{polylog Indep: Introduction}

Declarations of the Lean module \texttt{TCSlib.BooleanAnalysis.polylogIndep}.

\subsection*{polylog Indep: Results}

\begin{definition}[Boolean circuits]
\lean{Circuit}
\statementsource{lmn-constant-depth-fourier}{pdf-63759bd909d5-p004-b007}
\statementsource{Analysis-of-Boolean-Functions-by-Ryan-ODonnell}{pdf-ffbe9c7b3306-p103-b002}
\label{Circuit}
\leanok
A Boolean circuit on $n$ inputs is built from input variables indexed by
$\{0,\ldots,n-1\}$, Boolean constants, negations of circuits, and conjunctions or
disjunctions of finite lists of circuits.
\end{definition}

\begin{definition}[Circuit size]
\lean{Circuit.size}
\statementsource{lmn-constant-depth-fourier}{pdf-63759bd909d5-p004-b007}
\statementsource{Analysis-of-Boolean-Functions-by-Ryan-ODonnell}{pdf-ffbe9c7b3306-p104-b001}
\label{Circuit.size}
\leanok
\uses{Circuit}
The size of a circuit is defined recursively: an input or constant has size $1$;
negating a circuit adds $1$ to its size; and a conjunction or disjunction of a finite
list of circuits has size $1$ plus the sum of their sizes.
\end{definition}

\begin{definition}[Circuit depth]
\lean{Circuit.depth}
\statementsource{lmn-constant-depth-fourier}{pdf-63759bd909d5-p004-b007}
\statementsource{Analysis-of-Boolean-Functions-by-Ryan-ODonnell}{pdf-ffbe9c7b3306-p103-b002}
\label{Circuit.depth}
\leanok
\uses{Circuit}
The depth of a circuit is defined recursively: an input or constant has depth $0$, and
negating a circuit adds $1$ to its depth. A conjunction or disjunction of a finite list
has depth $1$ plus the maximum depth in the list, taking that maximum to be $0$ when the
list is empty.
\end{definition}

\begin{definition}[Circuit evaluation]
\lean{Circuit.eval}
\statementsource{Analysis-of-Boolean-Functions-by-Ryan-ODonnell}{pdf-ffbe9c7b3306-p103-b002}
\label{Circuit.eval}
\leanok
\uses{Circuit}
Given an assignment of Boolean values to the $n$ inputs, a circuit evaluates
recursively: an input returns its assigned value, a constant returns its Boolean value,
negation complements the subcircuit's value, conjunction returns true exactly when every
circuit in its list evaluates to true, and disjunction returns true exactly when at
least one circuit in its list evaluates to true.
\end{definition}

\begin{definition}[Limited independence]
\lean{IsKWiseIndependent}
\statementsource{lmn-constant-depth-fourier}{pdf-63759bd909d5-p012-b005}
\statementsource{Analysis-of-Boolean-Functions-by-Ryan-ODonnell}{
  pdf-ffbe9c7b3306-p146-b004,
  pdf-ffbe9c7b3306-p170-b003
}
\label{IsKWiseIndependent}
\leanok
Let $\mu$ be a probability distribution on Boolean assignments to $n$ variables, and let
$k$ be a nonnegative integer. The distribution is $k$-wise independent if, for every set
$S$ of variables with $|S|\le k$, the restriction of a random assignment drawn from
$\mu$ to $S$ is uniformly distributed over all Boolean assignments to $S$.
\end{definition}

\begin{definition}[Circuit acceptance probability]
\lean{Circuit.prob}
\label{Circuit.prob}
\leanok
\uses{Circuit, Circuit.eval}
For a circuit $C$ on $n$ Boolean inputs and a probability mass function $\mu$ on Boolean
assignments to those inputs, its acceptance probability is the sum, over all assignments
$x$, of $\mu(x)$ for those assignments on which $C$ evaluates to true.
\end{definition}

\begin{definition}[Circuit fooling condition]
\lean{Fools}
\statementsource{Analysis-of-Boolean-Functions-by-Ryan-ODonnell}{
  pdf-ffbe9c7b3306-p160-b003,
  pdf-ffbe9c7b3306-p167-b004
}
\label{Fools}
\leanok
\uses{Circuit, Circuit.prob, IsKWiseIndependent}
Let $k$ be a nonnegative integer, let $\varepsilon\in\mathbb{R}$, and let $S$ be a set
of Boolean circuits on $n$ inputs. The set $S$ is fooled with error $\varepsilon$ by
$k$-wise independent distributions if, for every $k$-wise independent probability
distribution $\mu$ on Boolean assignments to the inputs and every $C\in S$, the absolute
difference between the probability that $C$ accepts under $\mu$ and its acceptance
probability under the uniform distribution on all input assignments is at most
$\varepsilon$.
\end{definition}

\begin{definition}[Boolean monomial]
\lean{Monomial}
\label{Monomial}
\leanok
For a finite set $S\subseteq\{0,\ldots,n-1\}$ and a Boolean vector $x\in\{0,1\}^n$,
define its monomial to be the product over $i\in S$ of the indicator that $x_i=1$. Thus
the value is $1$ exactly when every coordinate indexed by $S$ is $1$, and is $0$
otherwise.
\end{definition}

\begin{definition}[Polynomial of bounded degree]
\lean{IsPolynomial}
\label{IsPolynomial}
\leanok
\uses{Monomial}
A function $P:\{0,1\}^n\to\mathbb R$ is a polynomial of degree at most $d$ if there are
real coefficients indexed by subsets $S\subseteq\{0,\ldots,n-1\}$, all zero whenever
$|S|>d$, such that at every $x\in\{0,1\}^n$, $P(x)$ is the sum over all such subsets of
the coefficient for $S$ times the Boolean monomial indexed by $S$ evaluated at $x$.
\end{definition}

\begin{theorem}[K-wise independence fools a monomial]
\lean{KWise_Fools_Monomial}
\label{KWise_Fools_Monomial}
\leanok
\uses{IsKWiseIndependent, Monomial}
\difficulty{6}
Let $n,k\in\mathbb N$, and let $\mu$ be a probability distribution on $\{0,1\}^n$ such
that every set of at most $k$ coordinates has a uniform joint distribution. For any set
$S$ of at most $k$ coordinates, the expectation under $\mu$ of the monomial that is $1$
exactly when every coordinate in $S$ is $1$ equals its expectation under the uniform
distribution on $\{0,1\}^n$.
\end{theorem}

\begin{theorem}[K-wise independence fools low-degree polynomials]
\lean{KWise_Fools_Polynomial}
\label{KWise_Fools_Polynomial}
\leanok
\uses{IsKWiseIndependent, IsPolynomial, KWise_Fools_Monomial, Monomial}
\difficulty{4}
Let $n,k\in\mathbb N$, and let $\mu$ be a probability distribution on $\{0,1\}^n$ such
that every set of at most $k$ coordinates has a uniform joint distribution. If $P$ is a
multilinear polynomial in the coordinate indicators, of degree at most $k$, evaluated on
the Boolean cube, then its expectation under $\mu$ equals its expectation under the
uniform distribution on $\{0,1\}^n$.
\end{theorem}

\begin{definition}[Sandwiching polynomials]
\lean{SandwichingPolynomials}
\label{SandwichingPolynomials}
\leanok
\uses{IsPolynomial}
For a function $f:\{0,1\}^n\to\mathbb R$, a natural number $k$, and a real number
$\varepsilon$, there are polynomials $P_\ell,P_u:\{0,1\}^n\to\mathbb R$ of degree at
most $k$ such that $P_\ell(x)\le f(x)\le P_u(x)$ for every $x$. Their average pointwise
gap under the uniform distribution on $\{0,1\}^n$ is at most $\varepsilon$.
\end{definition}

\begin{theorem}[Sandwiching polynomials fool independent distributions]
\lean{Sandwiching_Implies_Fooling}
\label{Sandwiching_Implies_Fooling}
\leanok
\uses{IsKWiseIndependent, KWise_Fools_Polynomial, SandwichingPolynomials}
\difficulty{4}
Let $n,k$ be nonnegative integers, let $f:\{0,1\}^n\to\mathbb{R}$, and let
$\varepsilon\in\mathbb{R}$. Suppose there are degree-at-most-$k$ polynomials
$P_\ell,P_u$, expressed as linear combinations of the monomials $x\mapsto\prod_{i\in
S}x_i$ for subsets $S\subseteq\{1,\ldots,n\}$, such that $P_\ell(x)\le f(x)\le P_u(x)$
for every $x$, and the uniform average of $P_u-P_\ell$ is at most $\varepsilon$. For
every probability distribution $\mu$ on $\{0,1\}^n$ whose marginal on each set of at
most $k$ coordinates is uniform, the absolute difference between the expectations of $f$
under $\mu$ and under the uniform distribution is at most $\varepsilon$.
\end{theorem}

\begin{definition}[Braverman sandwiching property]
\lean{BravermanProperty}
\statementsource{lmn-constant-depth-fourier}{pdf-63759bd909d5-p012-b005}
\label{BravermanProperty}
\leanok
\uses{Circuit, Circuit.depth, Circuit.eval, Circuit.size, SandwichingPolynomials}
There is a real constant $c>0$ such that for all natural numbers $n,d,m$ and every real
$ε>0$, if $m\ge1$, then every Boolean circuit on $n$ inputs with depth at most
$d$ and size at most $m$ has sandwiching polynomials for its Boolean-valued function,
with error $ε$ and degree at most
$\left\lceil\left(\log(m/\epsilon)+2\right)^{cd}\right\rceil$.
\end{definition}

\begin{lemma}[Ceiling bound for the Braverman parameter]
\lean{Braverman_Bound_Helper}
\label{Braverman_Bound_Helper}
\leanok
\difficulty{5}
Let $m,d$ be natural numbers and let $\varepsilon,c\in\mathbb{R}$ satisfy $m\geq1$,
$0<\varepsilon\leq1$, and $c>0$. Set $B=\log(m/\varepsilon)+2$, $E=cd$, and $k=\lceil
B^E\rceil$. Then $k\leq B^{(c+1)d}$.
\end{lemma}

\begin{lemma}[Probability is at most one]
\lean{Circuit.prob_le_one}
\label{Circuit.prob_le_one}
\leanok
\uses{Circuit, Circuit.prob}
\difficulty{3}
For a circuit $C$ on $n$ Boolean inputs and a probability mass function $\mu$ on Boolean
assignments, the probability that $C$ evaluates to true under $\mu$ is at most $1$.
\end{lemma}

\begin{lemma}[Probability is nonnegative]
\lean{Circuit.prob_nonneg}
\label{Circuit.prob_nonneg}
\leanok
\uses{Circuit, Circuit.prob}
\difficulty{1}
For a circuit $C$ on $n$ Boolean inputs and a probability mass function $\mu$ on Boolean
assignments, the probability that $C$ evaluates to true under $\mu$ is nonnegative.
\end{lemma}

\begin{lemma}[Zero-wise independence fools circuits]
\lean{Fools_zero_of_eps_gt_one}
\label{Fools_zero_of_eps_gt_one}
\leanok
\uses{Circuit, Circuit.prob, Circuit.prob_le_one, Circuit.prob_nonneg, Fools}
\difficulty{4}
Let $n$ be a natural number, let $\varepsilon>1$, and let $S$ be any set of Boolean
circuits on $n$ inputs. Every probability distribution on Boolean inputs that is
independent on every set of at most zero coordinates fools $S$ with error at most
$\varepsilon$: for every such distribution $\mu$ and every $C\in S$, the absolute
difference between the acceptance probabilities of $C$ under $\mu$ and under the uniform
distribution is at most $\varepsilon$.
\end{lemma}

\begin{lemma}[Probability as a weighted sum]
\lean{Circuit.prob_eq_sum}
\label{Circuit.prob_eq_sum}
\leanok
\uses{Circuit, Circuit.eval, Circuit.prob}
\difficulty{3}
For a circuit $C$ on $n$ Boolean inputs and a probability mass function $\mu$ on Boolean
assignments, its probability under $\mu$ equals $\sum_x
\mathbf{1}_{\{C(x)=\mathrm{true}\}}\,\mu(x)$, where the indicator is $1$ when $C(x)$ is
true and $0$ otherwise.
\end{lemma}

\begin{theorem}[Braverman theorem]
\lean{Braverman_Theorem}
\statementsource{lmn-constant-depth-fourier}{
  pdf-63759bd909d5-p012-b006,
  pdf-63759bd909d5-p012-b007,
  pdf-63759bd909d5-p013-b002
}
\label{Braverman_Theorem}
\leanok
\uses{BravermanProperty, Braverman_Bound_Helper, Circuit, Circuit.depth, Circuit.prob_eq_sum, Circuit.size, Fools, Sandwiching_Implies_Fooling}
\difficulty{4}
% dataset-pipeline: stub -- prose pending build_informal_statements.py
Statement pending regeneration.
\end{theorem}
